\begin{table}[ht]
\centering
\small
\caption{How many frequencies describe the signal. Residual, in units of the standard deviation of $\langle\Op\rangle(t)$, of the best single cosine at the splitting $\Delta E$ of the gated pair (what imaginary-time Monte Carlo of the splitting would predict) and at a free frequency; and the number of cosines of a free fit (greedy pursuit with joint refit of the amplitudes) needed for a residual below $20\%$, $10\%$ and $5\%$ ($>8$: not reached with up to eight cosines). Where the doublet test passes ($N=2,3$ and $N=4$, $k=20$) a splitting-based prediction is adequate; elsewhere the gated splitting explains none of the signal, and the number of cosines tends to grow with $N$ and to fall with $k$ (the basis is not converged at $N=5,6$). At $N=5$, $k=20$ the best free cosine sits at $\num{cr_N5_k20_wratio}$ times $\Delta E$.}
\label{tab:mc}
\begin{tabular}{lcccccc}
\toprule
 & \multicolumn{2}{c}{residual of one cosine} & \multicolumn{3}{c}{cosines for residual below} \\
\cmidrule(lr){2-3}\cmidrule(lr){4-6}
case & at $\Delta E$ & free & $20\%$ & $10\%$ & $5\%$ \\
\midrule
$N=2$, all $k$ & 0.10--0.14 & 0.10--0.13 & 1 & 1--2 & 3 \\
$N=3$, 6 basis sizes & 0.12--0.15 & 0.12--0.15 & 1 & 2--3 & 3--5 \\
$N=4$, $k=12$ & 1.00 & 0.72 & 4 & 7 & \ensuremath{>}8 \\
$N=4$, $k=16$ & 1.00 & 0.41 & 3 & 5 & 7 \\
$N=4$, $k=20$ & 0.18 & 0.17 & 1 & 3 & 5 \\
$N=5$, $k=12$ & 0.85 & 0.74 & 3 & 8 & \ensuremath{>}8 \\
$N=5$, $k=16$ & 1.00 & 0.48 & 3 & 4 & \ensuremath{>}8 \\
$N=5$, $k=20$ & 1.00 & 0.24 & 2 & 4 & 8 \\
$N=6$, $k=12$ & 0.99 & 0.55 & 4 & \ensuremath{>}8 & \ensuremath{>}8 \\
\bottomrule
\end{tabular}
\end{table}
